% =========================================================================================
% Chapter 4 -- Data Understanding (CRISP-DM)
%
% All counts here were read from the manifests, not from notes. Regenerate with
% scratchpad/dataset_facts.py and dataset_facts2.py if the manifests change.
% =========================================================================================
\section{Data Understanding}
\label{sec:data}

\subsection{Introduction}

This chapter describes the data the four perception modules were built on, and the single
question that had to be answered before any of them could be trained honestly: what counts
as one independent observation.

The question sounds procedural. It is not. Answering it wrongly inflates every metric that
follows, and does so invisibly --- the model appears to work, the validation score confirms
it, and nothing in the training process signals a problem. In each of the three imaging
datasets used here, the answer was not the one the file listing suggested.

\subsection{Dataset Overview}
\label{subsec:data-overview}

Four datasets were used, one per module. Tab.~\ref{tab:datasets} gives their composition.

\begin{table}[H]
  \centering
  \small
  \setlength{\tabcolsep}{5pt}
  \begin{tabular}{llrrl}
    \toprule
    \textbf{Module} & \textbf{Source} & \textbf{Files} & \textbf{Cases} &
    \textbf{Label type} \\
    \midrule
    Triage      & Four pooled ED datasets            & 164{,}700 & 164{,}700 &
      urgency tier (3) \\
    Cardiac     & CAMUS \cite{leclerc2019camus}      & 2{,}000 & 500 &
      pixel masks (3 structures) \\
    Gallbladder & UIdataGB (Mendeley)                & 10{,}694 & 220 &
      pathology class (9) \\
    Lung        & POCUS \cite{born2021pocus}         & 187 & 165 &
      findings (4, multi-label) \\
    \bottomrule
  \end{tabular}
  \caption{Composition of the four datasets. For the imaging datasets, ``files'' counts
    images, annotated frames or clips, and ``cases'' counts independent patients or
    acquisitions. The gap between the two columns is the subject of
    Section~\ref{subsec:data-grouping}.}
  \label{tab:datasets}
\end{table}

All four sources are publicly available and were obtained free of charge for academic
research. Their terms are not uniform, and the differences shaped how the project stores data.

The lung dataset \cite{born2021pocus} is a collection of contributions rather than a single
release, and the licence varies by contributing source: material from Neuruppin is CC~BY~4.0,
material from Northumbria is CC~BY--NC~4.0, and some items the original study used could not
be redistributed in its own repository at all. The cardiac dataset (CAMUS
\cite{leclerc2019camus}) and the gallbladder dataset (UIdataGB, Mendeley Data) are released
for research use. The triage sources are public-use files: NHAMCS is a United States federal
public-use release, the Iran ED data accompanies an open-access publication, and MC-MED and
MIMIC-IV-ED were used in their public demonstration releases, the credentialed full versions
being outside the scope of this work.

Two consequences follow, and both are visible in the deliverables. No source permits
wholesale redistribution of the media, so this project's repository contains
\emph{manifests} --- the lists of rows used, with their exclusion flags --- rather than
images, and reproduction requires obtaining the sources from their original providers. And
because several sources are non-commercial, the work reported here is academic; any use
beyond that would require re-examining each licence individually rather than treating the
collection as one permission.

\subsection{The Statistical Unit}
\label{subsec:data-grouping}

\begin{keystatement}
The gallbladder dataset contains 10{,}694 images. It does not contain 10{,}694 patients. It
contains 220.
\end{keystatement}

Medical imaging datasets are distributed as collections of files, and it is natural to treat
each file as one observation. For these datasets that is wrong, because the files are not
independent of one another. A single gallbladder examination contributed a median of 48
images --- and in one case 228 --- taken seconds apart, of the same organ, in the same
patient, under the same machine settings. Two such images are close to duplicates.

The consequence is specific. If a partition into training and validation data is drawn over
files, images from one examination land on both sides of it. The model then encounters, at
validation time, near-copies of images it memorised during training. It scores well. The
score measures recall of the training set, not generalisation to a new patient, and no
diagnostic within the training process distinguishes the two.

Fig.~\ref{fig:filesvscases} shows the size of the effect across the three imaging datasets.

\begin{figure}[H]
  \centering
  \includegraphics[width=0.86\textwidth]{fig_files_vs_cases.png}
  \caption{Files against independent cases for the three imaging datasets. The ratio bounds
    how far a file-level evaluation can overstate performance: the gallbladder dataset
    offers roughly 48 near-duplicate images per case, the cardiac dataset four annotated
    frames per patient, and the lung dataset only marginally more clips than patients.}
  \label{fig:filesvscases}
\end{figure}

The three datasets differ in how exposed they are. The lung dataset is nearly safe at
\(\times 1.1\); the cardiac dataset is moderately exposed at \(\times 4\); the gallbladder
dataset, at \(\times 48.6\), is where a file-level split would have produced a result with
no meaning at all.

\subsection{Recovering Case Identity}
\label{subsec:data-cases}

Knowing that the case is the unit is only useful if cases can be identified. The three
datasets required three different amounts of work.

\paragraph{Cardiac --- documented.} CAMUS is organised by patient. Each of the 500 patients
contributes two views (two-chamber and four-chamber), each annotated at end-diastole and
end-systole, giving 2{,}000 annotated frames. The grouping is read directly from the
directory structure and needs no reconstruction.

\paragraph{Gallbladder --- recoverable.} Case identity is not a column in the release, but
it is encoded in the filenames, which carry a per-case prefix. Parsing it yields 220 cases
across nine classes, with between 6 and 228 images each.

\paragraph{Lung --- reconstructed, with residual uncertainty.} The POCUS release carries no
patient identifier at all, while containing several clips of the same patient. Case identity
had to be inferred from filename patterns: explicit case numbers where present, day-series
markers for repeat examinations of one patient, and clip-numbering suffixes otherwise, with
a list of generic tokens excluded so that words such as \textit{video} or \textit{butterfly}
are not mistaken for patient identifiers. This recovered 165 cases from 187 files, with six
cases contributing more than one file and the largest contributing eight.

\begin{keystatement}
This last mapping is a heuristic, not ground truth. It is reported as such.
\end{keystatement}

Two clips of the same patient under different filename conventions may still be counted as
two cases, which would leave a small amount of leakage in place. The direction of that error
is known --- it can only inflate the lung results, never deflate them --- and it is carried
forward as a stated limitation rather than presented as a solved problem.
